\section{A flow-based Markov chain Monte Carlo algorithm}
\label{sec:flowMCMC}

\begin{figure*}
  \centering
  \begin{minipage}{.76\linewidth}
  \centering
  \subfloat[Normalizing flow between prior and output distributions]{
  \includegraphics[width=.97\linewidth]{images/FlowsDiagLabelled_taller.pdf}
  \label{subfl:norm-flow}
  }
  \end{minipage}%
  \begin{minipage}{.23\linewidth}
  \centering
  \subfloat[Inverse coupling layer]{
  \includegraphics[height=5.5cm]{images/coupling_layer.pdf}
  \label{subfl:inverse-coupling}
  }
  \end{minipage}
  \caption{ In \protect\subref{subfl:norm-flow}, a normalizing flow is shown transforming samples $z$ from a prior distribution $r(z)$ to samples $\phi$ distributed according to $\netp(\phi)$. The mapping $f^{-1}(z)$ is constructed by composing inverse coupling layers $g_i^{-1}$ as defined in Eq.~\eqref{eq:inv-affine-coupling} in terms of neural networks $s_i$ and $t_i$ and shown diagrammatically in \protect\subref{subfl:inverse-coupling}. By optimizing the neural networks within each coupling layer, $\netp(\phi)$ can be made to approximate a distribution of interest, $p(\phi)$.}
  \label{fig:flowdiag}
\end{figure*}

In lattice field theory, a Markov chain Monte Carlo (MCMC) process is an efficient way to generate field configurations $\phi \in \mathbb{R}^D$ distributed according to a target probability distribution
\begin{equation}
\begin{aligned}
    p(\phi) = e^{-S(\phi)} / Z, \quad
    \text{with}\quad Z =\int \prod_{j=1}^D d\phi_j \,e^{-S(\phi)},
\end{aligned}
\end{equation}
where $j$ indexes the $D$ components of $\phi$, $S(\phi)$ is the action that defines the theory and $Z$ is the partition function. Here, $\phi$ is defined to be a vector of $D$ real components representing the combined internal $(\alpha)$ and spacetime $(x)$ degrees of freedom of the field $\phi(x,\alpha)$ evaluated on a finite, discrete spacetime lattice (generalizations to gauge fields are discussed in Section~\ref{sec:summary}). A MCMC process generates a chain $\phi^{(0)} \rightarrow \phi^{(1)} \rightarrow \ldots \phi^{(N)}$ by steps through configuration space starting with an arbitrary configuration $\phi^{(0)}$. The steps are stochastic and are determined by the probabilities $T(\phi,\phi')$ associated with each possible transition $\phi \rightarrow \phi'$.
These probabilities must be non-negative and normalized:
%
\begin{equation}
    T(\phi,\phi')\ge 0
    \quad \text{and} \quad
    \int\prod_{j=1}^D d\phi'_j \, T(\phi,\phi') = 1.
\end{equation}
%
They must also satisfy the conditions of {\it ergodicity} and {\it balance} to ensure that samples in the chain are drawn from a distribution that converges to $p(\phi)$ after thermalization. For the chain to be ergodic, it must be possible to transition from a starting configuration $\phi$ to any other configuration $\phi'$ in a finite number of steps, i.e.,
%
\begin{equation}
    \Exists n \text{ such that } T^n(\phi,\phi')> 0 \text{ for all } \phi, \phi',
\end{equation}
%
%
and the chain must not have a period, for which it is sufficient that a single state has non-zero self-transition probability, i.e.,
%
\begin{equation}
    \Exists \phi \text{ such that } T(\phi,\phi) > 0.
\end{equation}
%
Balance is the condition that $p(\phi)$ is a stationary distribution of the transition:
%
\begin{equation}
    \int \prod_{j=1}^D d\phi_j \, p(\phi) T(\phi,\phi') = p(\phi').
\end{equation}
%
Any procedure which satisfies these conditions will, in the limit of a sufficiently long Markov chain, produce field configurations $\{\phi^{(i)}\}$ distributed according to $p(\phi)$.

\subsection{Metropolis-Hastings with generative models}
\label{subsec:generative-metropolis}
%
Given a model that allows sampling from a known probability distribution $\tilde{p}(\phi)$, a Markov chain for a desired probability distribution $p(\phi)$ can be constructed via the independence Metropolis sampler, a specialization of the Metropolis-Hastings method~\cite{tierney1994}. For each step $i$ of the chain, an {\it update proposal} $\phi'$ is generated by sampling from $\tilde{p}(\phi)$, independent of the previous configuration. This proposal is accepted with probability
%
\begin{equation}
A(\phi^{(i-1)},\phi') = \min\paren{1, \frac{\tilde{p}(\phi^{(i-1)}) }{ p(\phi^{(i-1)})} \frac{p(\phi')}{\tilde{p}(\phi')}}.
\label{eq:acc-rej}
\end{equation}
%
%
If the proposal is accepted, $\phi^{(i)} = \phi'$, otherwise $\phi^{(i)} = \phi^{(i-1)}$. This procedure defines the transition probabilities of the Markov chain.

The general Metropolis-Hastings algorithm has been proven to satisfy balance~\cite{hastings70} for any proposal scheme. For the independence Metropolis sampler, under the further condition that every state $\phi$ has non-zero proposal density and non-zero desired density,
%
\begin{equation}
    \tilde{p}(\phi) > 0, \; p(\phi) > 0 \text{ for all } \phi,
\end{equation}
%
the Markov chain is also ergodic and thus guaranteed to converge to the desired distribution~\cite{tierney1994}.

This Markov chain can be intuitively considered a method to correct an approximate distribution $\tilde{p}(\phi)$ to the desired distribution $p(\phi)$. The accept/reject statistics of the Metropolis-Hastings algorithm serve as a diagnostic for closeness of the approximate and desired distributions; if the distributions are equal, proposals are accepted with probability 1 and the Markov chain process is equivalent to a direct sampling of the desired distribution. This is made precise in Section~\ref{subsec:autocorr}.


\subsection{Sampling using normalizing flows}
\label{subsec:normalizing-flows}
Here, a {\it normalizing flow model} is used to define a proposal distribution $\tilde{p}(\phi)$ for a generative Metropolis-Hastings algorithm. Normalizing flows~\cite{rezende2015nf} are a machine learning approach to the task of sampling from complicated, intractable distributions. They do so by learning a map from an input distribution that is easy to sample to an output distribution that approximates the desired distribution. Normalizing flow models produce both samples and their associated probability densities, allowing the acceptance probability in Eq.~\eqref{eq:acc-rej} to be calculated.

A normalizing flow enacts the transformation between distributions by a change-of-variables\footnote{The convention of using $f^{-1}$ for the change-of-variables stems from typical applications of normalizing flows.}: a smooth, bijective function, $f^{-1}: \mathbb{R}^D \rightarrow \mathbb{R}^D$ maps samples $z$ from a prior distribution $r(z)$ to $\phi = f^{-1}(z)$. This mapping defines an output distribution $\netp(\phi)$, by the change-of-variables formula
\begin{align}
\label{eq:flowlikelihood}
\netp(\phi) = r(f(\phi))\left| \det  \frac{\partial f(\phi)}{\partial \phi} \right|.
\end{align}
Typically, the prior distribution is a simple and analytically-understood distribution (e.g., a normal distribution). While the desired distribution $p(\phi)$ is often complicated and difficult to sample from directly, optimizing the function $f$ allows one to generate samples from $\netp(\phi) \approx p(\phi)$. The function $f$ is chosen to have a tractable Jacobian such that the probability density $\netp(\phi)$ can be computed exactly according to Eq.~\eqref{eq:flowlikelihood}.
%




%
To encode a map from a simple distribution $r(z)$ to a complicated distribution $\netp(\phi)$, the map $f$ must be highly expressive while also being invertible and having a computable Jacobian. Here, the {\it real non-volume-preserving (real NVP) flow}~\cite{Dinh2016} machine learning approach is used: $f$ is constructed by the composition of {\it affine coupling layers} that scale and offset half of the components of the input at a time; the choice of which components of the data are transformed is part of the layer definition. Splitting the $D$-dimensional vector $\phi$ into ($D/2$)-dimensional pieces $\phi_a$ and $\phi_b$ according to this choice, a single coupling layer $g_i$ transforms $\phi$ to $z = g_i (\phi)$ via
%
\begin{align}\label{eq:affinecoupling}
    g_i(\phi) := \begin{cases}
    z_a = \phi_a  \\
    z_b = \phi_b \odot e^{s_i(\phi_a)} + t_i(\phi_a),
    \end{cases}
\end{align}
%
where $s_i$ and $t_i$ are neural networks mapping from $\mathbb{R}^{D/2}$ to $\mathbb{R}^{D/2}$ and $\odot$ denotes element-wise multiplication. Importantly, each layer $g_i$ is invertible without inverting the neural networks $s_i$ or $t_i$:
%
\begin{align}\label{eq:inv-affine-coupling}
    g_i^{-1}(z) := \begin{cases}
    \phi_a = z_a  \\
    \phi_b = (z_b - t_i(z_a)) \odot e^{-s_i(z_a)}.
    \end{cases}
\end{align}
%
The Jacobian matrix is lower-triangular and its determinant can be easily computed. For coupling layer $g_i$:
%
\begin{equation}
    \left| \det \pbp{g_i(\phi)}{\phi} \right| = \prod_{j=1}^{D/2} e^{\sq{s_i(\phi_a)}_j},
\end{equation}
where $j$ indexes the $D/2$ components of the output of $s_i$.
%
Stacking many coupling layers $g_1, \dots, g_n$ which alternate which half of the data is transformed, the function $f$ is defined as
%
\begin{equation}
    f(\phi) = g_1 (g_2 ( \dots g_{n} (\phi) \dots ) ).
\end{equation}
%
Using the chain rule, the determinant of the Jacobian of $f$ is a product of the contributions from each $g_i$. By increasing the number of coupling layers and the complexity of the networks $s_i$ and $t_i$, $f$ can systematically be made more expressive and general. Figure~\ref{fig:flowdiag} depicts how composing many coupling layers incrementally modifies a prior distribution which is easy to sample into a more complex output distribution that approximates a distribution of interest.

For a fixed initial distribution $r(z)$, the neural networks within each affine coupling layer of $f$ can be trained to bring $\netp(\phi)$ close to the desired distribution $p(\phi)$. This training is undertaken by minimizing a {\it loss function}. Here, the loss function used is a shifted Kullback-Leibler (KL) divergence\footnote{This training paradigm is a specific instance of Probability Density Distillation~\cite{Oord2017WaveNet, Li2018}.
} between the target distribution of the form $p(\phi) = e^{-S(\phi)}/Z$ and the proposal distribution $\netp(\phi)$:
%
\begin{equation}
\begin{aligned}\label{eq:loss}
L(\netp) &:= D_{KL}(\netp || p) - \log Z \\
&= \int \prod_j d\phi_j \, \netp(\phi) \paren{\log\netp(\phi) - \log p(\phi) - \log Z} \\
&=  \int \prod_j d\phi_j \, \netp(\phi) \paren{\log\netp(\phi) + S(\phi)}.
\end{aligned}
\end{equation}
%
This loss function has been successfully applied in related generative approaches to statistical lattice models~\cite{DBLP:journals/corr/abs-1809-10188,Li2018}. The formal shift by $\log{Z}$ in Eq.~\eqref{eq:loss} eliminates the need to compute the true partition function, and does not affect the gradients or location of the minima. By non-negativity of the KL divergence, the lower bound on the loss is $-\log{Z}$, and this minimum is achieved exactly when $\netp = p$.
In practice, the loss is stochastically estimated by drawing batches of $M$ samples from the model $\curly{\phi^{(i)} \sim \netp}$ and computing the sample mean:
%
\begin{equation}
\begin{aligned}\label{eq:loss-batch}
    \widehat{L\paren{\netp}} =
    \frac{1}{M} \sum_{i=1}^M \paren{\log \netp(\phi^{(i)}) + S(\phi^{(i)})}.
\end{aligned}
\end{equation}
%
The loss minimization can then be undertaken using stochastic optimization techniques such as stochastic gradient descent or momentum-based methods including Adam and Nesterov~\cite{Kingma2015Adam,nesterov1983}.

By construction, the flow model allows sampling from $\netp$ efficiently. The training process can thus be performed by drawing samples from the model itself, rather than using existing samples from the desired distribution as training data. This self-training is a key feature of the proposed approach to Monte Carlo sampling for field theories, where samples from the desired distribution are often computationally expensive to obtain. If samples do exist, they can be used to `pre-train' the network, although in the numerical studies undertaken here this was not found to be markedly more efficient in network optimization than using only self-training.

Given a trained model with distribution $\netp(\phi) \approx p(\phi)$, samples from $\netp$ can be used as proposals to advance a Markov chain using the generative Metropolis-Hastings algorithm described above. This forms the basis for the \emph{flow-based MCMC} algorithm proposed here:
\begin{enumerate}
    \item A flow-based generative model (here, a {real NVP} model) is trained using the shifted KL loss given in Eq.~\eqref{eq:loss} to have output distribution $\netp(\phi) \approx p(\phi)$;
    \item $N$ proposals $\curly{\phi'^{(i)} \sim \netp}$ are produced by sampling from the flow-based model (this can be done in parallel) and the associated action $S(\phi')$ is computed for each proposal;
    \item Starting from an arbitrary initial configuration, each proposed sample is successively accepted or rejected using the Metropolis-Hastings algorithm given in Eq.~\eqref{eq:acc-rej} to build a Markov chain of length $N$.
\end{enumerate}
%
When the prior distribution $r(z)$ is strictly positive, the invertibility and continuity of $f$ guarantees that the generated distribution $\netp(\phi)$ is also strictly positive. For all models with finite action, and thus $p(\phi) > 0$, the resulting Markov chain is then ergodic by the arguments detailed in Section~\ref{subsec:generative-metropolis}.


\subsection{Autocorrelation time for generative Metropolis-Hastings}
\label{subsec:autocorr}

For any Markov chain constructed via a generative Metropolis-Hastings algorithm (with independent update proposals), an observable-independent estimator for autocorrelation time can be defined from the accept/reject statistics of the chain. This serves both as a measure of the similarity between the proposal and desired distributions and enables proper error estimation for lattice observables~\cite{Wolff2003}.

Precisely, the autocorrelation at Markov chain separation $\tau$, for all observables, is given by the probability of $\tau$ rejections in a row,
%
\begin{equation}
p_{\tau\text{rej}} \equiv \avg{\prod_{i = 1}^\tau \mathbbm{1}_{\text{rej}}(i)} = \rho(\tau)/\rho(0),
\end{equation}
%
where $\mathbbm{1}_{\text{rej}}(i)$ is an indicator variable taking value $1$ when the proposed step from $i-1$ to $i$ was rejected in the Metropolis-Hastings algorithm, and $0$ otherwise. In practice, for a near-equilibrium, finite Markov chain with length $N$, a finite-sample estimator provides a good approximation to $\rho(\tau)/\rho(0)$:
%
\begin{equation}
\widehat{\rho(\tau)/\rho(0)}_{\text{acc}} = \frac{1}{N-\tau} \sum_{j=1}^{N-\tau} \prod_{i = 1}^\tau \mathbbm{1}_{\text{rej}}(i+j).
\end{equation}
%
This measure of autocorrelation is consistent with the usual definition; it is shown in Appendix~\ref{app:bounds} that the standard estimator for the autocorrelation of any given observable $\obs$,
%
\begin{equation}
\widehat{\rho(\tau)/\rho(0)}_\obs
=  \frac{\frac{1}{N-\tau} \sum_{i=0}^{N-\tau-1} (\obs_i - \bar{\obs})(\obs_{i+\tau}-\bar{\obs})}{\frac{1}{N} \sum_{i=0}^{N-1} \paren{\obs_i - \bar{\obs}}^2},
\end{equation}
%
also converges to $p_{\tau\text{rej}}$ in the limit $N\rightarrow \infty$.

The qualitative relation between acceptance rate and autocorrelations gives a convenient measure of the autocorrelation characteristics of a Markov chain. Precisely, the autocorrelation at distance $\tau$ can be bounded in terms of the average acceptance rate $a = 1 - \mathbb{E}\sq{p_\text{rej}}$:
\begin{equation}
\rho(\tau)/\rho(0) = \expt{\phi \sim p}\sq{p_{\text{rej}}^\tau(\phi)} \geq \paren{\expt{\phi \sim p}\sq{p_{\text{rej}}(\phi)}}^\tau = (1-a)^\tau.
\label{eq:mean-acc-bound}
\end{equation}
Increasing the acceptance rate of an independence Metropolis sampler is thus a necessary condition to reduce autocorrelations. Additionally, $a = 1$ exactly when the proposal and desired distributions are equal. In this case, $p_\text{rej}(\phi) = 0$ for each $\phi$, and there are no autocorrelations. While bringing $a$ close to $1$ does not provide an upper bound on autocorrelation, stochastically improving a loss function that measures distance between distributions is expected to reduce autocorrelations on average. In practice, autocorrelations should be evaluated as a test metric alongside the training loss to confirm improvement over the course of training the model. The correspondence between loss minimization, acceptance rate, and autocorrelations is studied in the context of $\phi^4$ theory in Section~\ref{sec:phi4}, where a clear correlation between $a$ and $\rho(\tau)/\rho(0)$ is observed.

\subsection{Critical slowing down}
\label{subsec:CSD-theory}

When a distribution is sampled using a Markov chain with large autocorrelation time, many updates are required to produce decorrelated samples. Critical slowing down (CSD) is defined as the divergence of the autocorrelation time of Markov chain sampling as a critical point in parameter space is approached~\cite{Wolff:1989wq}.
%
A numerically-stable definition of the characteristic autocorrelation time of a Markov chain is the \emph{integrated autocorrelation time}:
%
\begin{equation}
    \tau_\mathcal{O}^\text{int} = \frac{1}{2} + \lim_{\tau_\text{max} \rightarrow \infty}\sum_{\tau=1}^{\tau_\text{max}} \frac{\rho_\mathcal{O}(\tau)}{\rho_\mathcal{O}(0)}.
\end{equation}
As a critical point is approached, analysis of standard local-update algorithms for lattice models suggests $\tau_\mathcal{O}^\text{int}$ is typically well-described by a power law in the lattice spacing, or for fixed physical volume, a power law in the lattice sites per dimension, $L$. A dynamical critical exponent $z_\mathcal{O}$ is thus defined by a fit to $\tau_\mathcal{O}^\text{int} = \alpha_\mathcal{O} L^{z_\mathcal{O}}$ along a line of constant physics. An update algorithm for which the critical exponent is zero is unaffected by CSD.

In any generative Metropolis-Hastings simulation, the autocorrelation time is completely fixed by the expected accept/reject statistics, which in turn result from the structure of the proposal and desired distributions. For models trained with a target value of the integrated autocorrelation time used as a stopping criterion, CSD associated with the Markov chain sampling is thus trivially removed at the expense of up-front training costs. The difficulty of CSD is in essence shifted to the training of the model, i.e., to the optimization of the proposal distribution. The cost and scaling of this optimization task for $\phi^4$ theory, as studied here, and the prospects of scaling this approach to more complicated theories, are discussed in Sec.~\ref{subsec:training-cost}.
